%%!TEX root = ../../template.tex
\typeout{NT FILE chapter5}

\chapter{System Evaluation}
\label{cha:evaluation}

\prependtographicspath{{Figures/}}

In this chapter we present the experimental evaluation of \gls{SEnTriS}. We start by analyzing and discussing the solution's storage and query complexity, in Section \ref{ch5:complexity-analysis}. In Section \ref{ch5:setup-methodology}, we present the experimental evaluation conducted: we describe the testing environments; explain the experimental evaluation scenarios; present and discuss the results. Finally, in Section \ref{ch5:security-analysis}, we present \gls{SEnTriS} security properties, analyze the leakage, and discuss the trade-offs between each version of the protocol.

\section{Complexity Analysis}
\label{ch5:complexity-analysis}

In this Section we will evaluation \gls{SEnTriS} in terms of space complexity for storage of the encrypted data, and time complexity for query evaluation.

\subsection{Storage Complexity}

When evaluating the storage complexity of \gls{SEnTriS}, we need to take into consideration that two different types of encryption operations are done, which generate ciphertexts with different sizes. The size of a ciphertext generated with the XChaCha20-Poly1305 algorithm, depends on the size of the plaintext (plus a constant 16 bytes tag), whereas the size of the ciphertexts generated with the \gls{DGK} cryptosystem depend on the size of the modulus $n$ (usually with 1024 or 2048 bits).

\paragraph{SEnTriS V1 (\emph{Abox})} For every triple parsed, we generate 9 trapdoors (one per keyword combination) and their respective encrypted node. Furthermore, we also generate the trapdoors pointing to the encrypted frequencies of the triple constituents. Finally, for each keyword we generate a trapdoor referencing their encrypted frequencies. 

\paragraph{SEnTriS V2 (\emph{Abox})} For every triple parsed, we generate 9 trapdoors (one per keyword combination), their respective encrypted node and the nodes equality tag, which is also encrypted and stored twice. Furthermore, we generate the trapdoors pointing to the encrypted frequencies of the triple constituents. Finally, for each keyword we generate a trapdoor referencing their encrypted frequencies, and for each node we generate a trapdoor pointing to its encrypted equality tag.

\paragraph{SEnTriS V1 \& V2 (\emph{Tbox})} For every triple, we generate 3 trapdoors referencing each of its nodes, which are also encrypted.

Let $T_a$, $K_a$ and $N_a$ be, respectively, the triples, keywords and nodes of a dataset's \emph{Abox}. Let $T_t$, $K_t$ and $N_t$ be, respectively, the triples, keywords and nodes of the same dataset's \emph{Tbox}. In Table \ref{tbl:storage-complexity:count} we summarize the count of encryption functions executed during an upload (considering no deletions).

%Space complexity in terms of RDF terms

\begin{table}[!ht]
    \centering
    \begin{tabular}{|l|l|l|l|}
    \hline
         \textbf{Version} & \textbf{Deterministic} & \textbf{Probabilistic} & \textbf{DGK} \\ \hline
        V1 & $27T_a + K_a + 3T_t$ & $18T_a + Ka + 3T_t$ & 0 \\ \hline
        V2 & $18T_a + K_a + N_a + 3T_t$ & $18T_a + K_a + N_a + 3T_t$ & $18T_a$ \\ \hline
    \end{tabular}
    \caption{Count of encryption operations.}
    \label{tbl:storage-complexity:count}
\end{table}

Now, let $\overline{S}$ be the average size of a ciphertext generated with the XChaCha20-Poly1305 and $\overline{N}$ the size of a ciphertext generated with the \gls{DGK} cryptosystem. In Table \ref{tbl:storage-complexity:formulas} we show the storage complexity for each of the \gls{SEnTriS} versions.

\begin{table}[h!]
    \centering
    \begin{tabular}{|l|l|l|}
    \hline
        \textbf{Version} & \textbf{Storage Formula} & \textbf{Asymptotic Storage} \\
        \hline
        V1 & $(45T_a + 2K_a + 6T_t) * \overline{S}$ & $O\!\left((T_a + K_a + T_t) * \overline{S}\right)$ \\
        \hline
        V2 & $(36T_a + 2K_a + 2N_a + 6T_t) * \overline{S} + 18T_a * \overline{N}$  & $O\!\left((T_a + K_a + N_a + T_t) * \overline{S} + T_a * \overline{N}\right)$ \\
        \hline
    \end{tabular}
    \caption{Storage complexity comparison.}
    \label{tbl:storage-complexity:formulas}
\end{table}


To conclude, in both versions storage complexity scales linearly. Although, in \gls{SEnTriS} V2 the number of ciphertexts increases by approximately $9T_a + 2N_a$ relative to \gls{SEnTriS} V1, the storage overhead is substantially higher due to the increase in size of the ciphertexts generated with the \gls{DGK} cryptosystem.

\subsection{Query Evaluation Complexity}
%Time complexity in terms of node search, number of joins and modular exponentiations. Also consider transitivity and expansion depth
\subsubsection{SEnTriS V1}
\subsubsection{SEnTriS V2}

\section{Experimental Evaluation}
\label{ch5:setup-methodology}

\subsection{Setup \& Methodology} 
\subsubsection{Validation Scenarios}
\subsubsection{LUBM Scenarios}
\subsection{Experimental Results}
\subsubsection{Repository Size \& (Up)load Time}
\subsubsection{Query Response Time}
\subsubsection{Query Completeness \& Soundness}
\subsubsection{Result Analysis}

\section{Security Analysis}
\label{ch5:security-analysis}
In this Section we will discuss the security of \gls{SEnTriS}, taking into consideration its security model, defined in Section \ref{ch3:sec:solution-models:security-model}. We will start by enumerating the security properties of the different versions.

\subsection{Security Properties}

We expect all implementations of \gls{SEnTriS} to possess the following \emph{security properties}:

\begin{enumerate}
    \item A user will be able to create triplestores and access tokens, if and only if they have an active session.
    \item A user will be able to interact with a triplestore, if and only if they have a valid access token.
    \item A user will not be able to query triplestores which they have not been given access to.
    \item A user will only be able to execute operations according to its role.
    \item An external observer will not be able to learn any useful information from eavesdropping messages sent between components.
    \item An attacker will not be able to impersonate any component.
    \item An attacker will not be able to tamper with messages sent between components.
    \item An attacker with access to the \emph{Triplestore} component will learn nothing more than the size of a triplestore and its \emph{search} and \emph{access} patterns.
    \item Regarding \emph{term equality} during intermediate set operations between \emph{bindings}, an attacker with access to the \emph{Proxy} component will learn nothing more than:
        \begin{enumerate}
            \item \emph{SEnTriS V1}: the equality of all terms, belonging to a series of queries, targeting the same triplestore.
            \item \emph{SEnTriS V2} the equality of all terms present in the graph patterns of individual queries. The bindings of different queries are indistinguishable.
        \end{enumerate}
\end{enumerate}

% Plan V0 [[search_0,0, ... search0,n]->[random join_pipeline]^K -> [UNION|OPTIONAL|MINUS|JOIN search0,0, join_0,n, join,k,n]-> DISTINCT | PROJECT | ORDERBY  | SLICE]

\subsection{Leakage Analysis}
\subsubsection{SEnTriS V1}
\subsubsection{SEnTriS V2}
\subsection{Security Properties Proofs}

%Corolary: sentry v2 leakage is reduced

\section{Final Discussion}

Both implementations come with a trade-off between query evaluation leakage and performance.

\gls{SEnTriS} V1 provides both data confidentiality and access control, but suffers from the same leakage problems of \emph{CryptDB}, where an adversary can associate the leakage produced by multiple queries targeting the same triplestore. This facilitates the computation of the frequency distribution of all observed terms, which simplifies the execution of \emph{leakage-abuse} attacks \autocite{Jutla2021-af, Hahn2019-sg, Shafieinejad2021-rb}. This can be especially problematic during \gls{SPARQL} queries evaluation, as they tend to require a high volume of equality checks. A situation further aggravated when we apply query rewriting techniques: the modified queries usually require an even higher number of equality joins.

In \gls{SEnTriS} V2 provides both data confidentiality, access control and reduced leakage. This is achieved at the cost of: 1) an increase in complexity, during set operations between binding tables: 2) an extra communication round during query evaluation, where we have to reverse the masking applied to the \emph{equality tags} and need to fetch, and decrypt, the node values to generate the final results.